This paper presented a cascaded spatio-temporal framework for GI abnormality detection in endoscopy that combines a
MobileNetV2 gate, a ResNet-50 feature extractor and a Bi-LSTM with temporal attention, trained with focal loss and
explained with Grad-CAM. It was benchmarked against a classical SIFT-BoVW+SVM pipeline and eight deep baselines on
four public databases with classification, efficiency and explainability evaluators. On clean still images the
proposed system comes within \Rgapbest\ points of the best single-frame network (\Rpropacc\% accuracy and \Rpropmcc\%
MCC on Kvasir v2, versus \Rbestbaseacc\% for \Rbestbase). Its
advantage appears when frames are degraded: under transient motion blur, defocus, specular flashes or occlusion it
keeps \Rrobtransient\% accuracy, whereas single-frame baselines fall to \Rrobbestbaseacc\% or less. A temporally
smoothed gate retains $\ge$99\% of abnormal frames while skipping the expensive stages for normal ones, which cuts
the average compute per frame and pays off increasingly as the share of normal frames grows, as in real examinations. Future work will
validate the system on real endoscopic video, add repeated-seed confidence intervals, and extend the label set to
rare gastric pathologies such as early gastric cancer and ulcers.
